% TOP_PROBLEM: {"id": 30006689, "title": "Prove the Penrose inequality or present a counterexample", "category_id": 16, "category": {"id": 16, "name": "physics", "display_name": "Mathematical Physics", "description": "Problems at the intersection of mathematics and physics.", "slug": "physics", "order_index": 16, "created_at": "2026-07-31T15:26:25.671Z"}, "status": "open"}
% FIRST_POSTED: 2026-09-27
% !TeX program = lualatex
% Standalone notebook; compile twice: lualatex 30006689.tex
% Original section preserved verbatim except for marked insertion blocks.
\documentclass[11pt,a4paper]{article}
\usepackage[margin=24mm,headheight=14pt]{geometry}
\usepackage{fontspec}
\setmainfont{Latin Modern Roman}
\setsansfont{Latin Modern Sans}
\setmonofont{Latin Modern Mono}
\usepackage{amsmath,amssymb,mathtools}
\usepackage{unicode-math}
\setmathfont{Latin Modern Math}
\usepackage{microtype}
\usepackage{enumitem,xurl,needspace}
\usepackage[dvipsnames]{xcolor}
\usepackage{hyperref}
\hypersetup{unicode=true,colorlinks=true,linkcolor=MidnightBlue,urlcolor=MidnightBlue,
 pdftitle={Research notebook: Problem 30006689},
 pdfauthor={Open Problems Math research notebook},bookmarksnumbered=true}
\usepackage{bookmark,titlesec,fancyhdr}
\titleformat{\section}{\Large\bfseries\raggedright}{\thesection}{0.7em}{}
\pagestyle{fancy}\fancyhf{}
\fancyhead[L]{\small\sffamily Open Problems Math\enspace|\enspace Research notebook}
\fancyhead[R]{\small\sffamily Problem 30006689}
\fancyfoot[L]{\small\sffamily 27 September 2026}
\fancyfoot[R]{\small\sffamily \thepage}
\setlength{\parindent}{0pt}
\setlength{\parskip}{5pt plus 1pt}
\setlength{\emergencystretch}{3em}
\setcounter{secnumdepth}{1}
\setcounter{tocdepth}{1}
\setlist[itemize]{leftmargin=3em,itemsep=5pt,topsep=4pt,parsep=0pt}
\allowdisplaybreaks
\newcommand{\N}{\mathbb{N}}
\newcommand{\Z}{\mathbb{Z}}
\newcommand{\Q}{\mathbb{Q}}
\newcommand{\R}{\mathbb{R}}
\newcommand{\C}{\mathbb{C}}
\newcommand{\F}{\mathbb{F}}
\newcommand{\A}{\mathbb{A}}
\newcommand{\PP}{\mathbb P}
\newcommand{\E}{\mathbb{E}}
\newcommand{\eps}{\varepsilon}
\newcommand{\dd}{\,\mathrm d}
\newcommand{\sref}[2]{\hyperlink{source:#1:#2}{[S#2]}}
\newcommand{\eref}[2]{\hyperlink{extra:#1:#2}{[E#2]}}
\newif\ifshowcatalogue
\showcataloguetrue
\newcommand{\cataloguescope}[1]{\ifshowcatalogue
 \paragraph{Exact scope in the supplied catalogue.}
 \begin{quote}\small #1\end{quote}\fi}
\numberwithin{equation}{section}
\begin{document}
\setcounter{section}{0}
% BEGIN PRESERVED SECTION
% ================================================================
% problemId: 30006689
% Canonical problem ID: 30006689
% exactTarget SHA-256: e996a87efab6ebd50e93f71948051eeb0bd438fce5af1d11a5edcaa4653db242
\clearpage
\section*{\texorpdfstring{Prove the Penrose inequality or present a counterexample}{Prove the Penrose inequality or present a counterexample}}\label{30006689}
\subsection{Definitions and mathematical statement}
An initial data set consists of a Riemannian $3$-manifold $(M,g)$ and symmetric tensor $k$, its second fundamental form. In geometric units the constraint densities obey
\[
 16\pi\mu=R_g+(\operatorname{tr}_gk)^2-|k|_g^2,\qquad
 8\pi J=\operatorname{div}_g(k-(\operatorname{tr}_gk)g).
\]
The dominant energy condition is $\mu\ge|J|_g$. At an asymptotically Euclidean end, ADM energy is defined by the limiting metric-flux integral; the mass convention and decay hypotheses are those of the cited formulation. If $A$ is the infimum of areas enclosing the specified outermost apparent horizon, the desired bound is
\[
 m_{\rm ADM}\ge\sqrt{A/(16\pi)},
\]
with the source's Schwarzschild rigidity condition. An apparent horizon is characterized by vanishing appropriate outgoing null expansion. The full spacetime statement does not assume $k=0$.
\par\smallskip\textit{Formulation and concept references:} \sref{30006689}{1}.
\cataloguescope{Let (M,g,k) be a three-dimensional asymptotically flat initial data set for the Einstein equations satisfying the dominant energy condition and containing an outermost apparent horizon. If A is the least area needed to enclose that horizon and m\_ADM is the ADM mass, prove the sharp spacetime Penrose inequality m\_ADM >= sqrt(A/(16*pi)), with the appropriate Schwarzschild rigidity statement, or construct initial data satisfying the hypotheses that violate the inequality.}
\subsection{Short English statement}
Does the mass of an asymptotically flat spacetime initial state always dominate the area-based lower bound from its outermost horizon?
% BEGIN RESEARCH 37
\subsection{Research attempt}

\paragraph{Current frontier.}
The selected problem is the sharp spacetime inequality, with the least
\emph{enclosing} area and Schwarzschild rigidity. The Riemannian $k=0$ case
and spherical cases must be kept separate. Bray--Khuri develop a generalized
Jang approach and prove its spherical existence theory. \eref{30006689}{1}.
Allen--Bryden--Kazaras--Khuri's Theorem~1.1 supplies a universal but suboptimal
constant; it cannot be changed to $16\pi$ by rescaling the data.
\sref{30006689}{3}. Ellithy's May 2026 manuscript retains an explicit quasi-final-state
hypothesis. \sref{30006689}{2}. The abstract, formulation passages, and relevant
analytic passages of Xu's 510-page v5 were inspected, not the entire proof.
Its extension from an area-maximizing MOTS to the outermost MOTS is explicitly
conditional. The recorded target is preserved, but neither that conditional
extension nor its different area functional is silently adopted as a theorem
about the target's least enclosing area. \sref{30006689}{1}.

\paragraph{Attack route.}
A spacetime monotonic mass should connect the horizon area to the ADM end.
I first derive the sharp mechanism in a spherical exterior without setting
$k=0$, then ask which term loses its sign outside symmetry. The alternative
Jang/conformal strategy must control both divergence terms and singular
limits. A source-audit test below isolates an endpoint estimate that cannot
justify such a limit on its own.

\paragraph{Attempt.}
Work in the explicitly restricted exterior $[s_0,\infty)\times S^2$ with
\[
 g=ds^2+r(s)^2d\Omega^2,\qquad
 k=\kappa(s)ds^2+\tau(s)r(s)^2d\Omega^2.
\]
Assume the source's asymptotically flat end conditions, $r\to\infty$,
$r'\to1$, $r^3\tau^2\to0$, and existence of the ADM flux limit. Put
$J_s=J(\partial_s)$. Direct evaluation of the constraints gives
\begin{align*}
16\pi\mu&=-4r''/r+2(1-r'^2)/r^2+4\kappa\tau+2\tau^2,\\
8\pi J_s&=2(r'/r)(\kappa-\tau)-2\tau'.
\end{align*}
For example $\operatorname{tr}k=\kappa+2\tau$ and
$|k|^2=\kappa^2+2\tau^2$, producing $4\kappa\tau+2\tau^2$.
Define the spherical spacetime Hawking mass
\[
 m(s)=\frac r2(1-r'^2+r^2\tau^2).
\]
Differentiating, substituting the two constraint identities, and canceling
$\kappa\tau$ yields the exact identity
\begin{equation}\label{r37:identity}
 m'=4\pi r^2(\mu r'-r\tau J_s).
\end{equation}
This is checked symbolically in the companion code as well as by the displayed
substitution. The formula is valid for either sign of $\tau$.

Assume a future apparent-horizon sphere at $s_0$, so in this sign convention
$r'(s_0)+r(s_0)\tau(s_0)=0$, and impose the additional untrapped-exterior
condition
\begin{equation}\label{r37:extra}
 r'>|r\tau|\quad(s>s_0).
\end{equation}
This is an explicit restriction, not a consequence claimed for general
outermost-MOTS data. The dominant energy condition gives
\[
 m'\ge4\pi r^2\mu(r'-|r\tau|)\ge0,
 \qquad m(s_0)=r(s_0)/2.
\]
In areal coordinates the metric is $dr^2/r'^2+r^2d\Omega^2$; its ADM energy
is $\lim r(1-r'^2)/2$ under the stated decay, and spherical symmetry gives
zero ADM momentum. The $\tau$ term tends to zero, hence
$m_{\rm ADM}=\lim m(s)\ge r(s_0)/2$.

This really controls the least enclosing area in this restricted exterior,
not merely the area of a preferred coordinate surface. Write $r_0=r(s_0)$.
The closed form $\omega=r_0^2\sin\theta\,d\theta\wedge d\phi$ has comass
$r_0^2/r(s)^2\le1$, since \eqref{r37:extra} makes $r$ increasing. Every smooth
oriented enclosing surface homologous to the horizon has
$\int_\Sigma\omega=4\pi r_0^2$ by Stokes, so
$|\Sigma|\ge4\pi r_0^2$. The horizon, or spheres approaching it from outside,
attains that infimum. Thus $A=4\pi r_0^2$ and
\[
 m_{\rm ADM}\ge\sqrt{A/(16\pi)}
\]
in this class. This recovers a known kind of spherical result by a direct
calculation; it is not claimed as a new theorem resolving general data.
If equality holds, \eqref{r37:extra} and continuity force $\mu=J_s=0$ outside
the horizon. The calculation itself establishes that necessary vacuum
condition, not a complete nonspherical Schwarzschild rigidity argument.

The sharp constant cannot be recovered from a suboptimal one by changing
units. Under the rescaling of all lengths by $c>0$, mass is multiplied by $c$
and area by $c^2$, leaving $m/\sqrt A$ unchanged. Likewise \eqref{r37:identity}
explains the real obstruction to extending this proof: an arbitrary exterior
has no globally defined areal-radius foliation for which the corresponding
momentum and geometric terms reduce to $\mu r'-r\tau J_s$ with
$r'>|r\tau|$. Existence of such a controlled foliation is additional analysis.

\textbf{Endpoint stress test for a different proposed route.}
Xu v5, pp.~342--343, states an $L^{3/2}\to L^\infty$ bound for the Newtonian
potential in its discussion of smoothing, before mentioning an upgraded
$L^2$ estimate. The endpoint assertion alone is false already in Euclidean
space. \sref{30006689}{1}. To check this independently, let
\[
 f_N(x)=\frac{1_{\{e^{-N}<|x|<e^{-1}\}}}
 {|x|^2\log(1/|x|)}\qquad(N>2).
\]
Then spherical integration gives
\[
 \|f_N\|_{3/2}^{3/2}=4\pi\int_1^N t^{-3/2}\,dt\le8\pi,
 \qquad \frac1{4\pi}\int\frac{f_N(x)}{|x|}\,dx=\log N.
\]
Thus there is no uniform $L^{3/2}\to L^\infty$ operator bound. The same
conclusion holds with smooth nonnegative cutoffs equal to one on a slightly
smaller annulus: the norm stays bounded and the value still grows like
$\log N$. Scaling the functions can even make their $L^{3/2}$ norms tend to
zero while the potential at zero diverges. This tests one analytic sentence,
not all subsequent estimates: an independently established stronger $L^2$
or Lorentz-space control could avoid it. No claim to have refuted the whole
510-page manuscript is made.

\paragraph{Stress test.}
The spherical proof used \eqref{r37:extra}; dropping it removes the sign of
\eqref{r37:identity}, even if the dominant energy condition remains true.
Choosing a larger non-outer-minimizing horizon area would also change the
problem. The source's equality convention refers to a slice of Schwarzschild
spacetime, not necessarily the time-symmetric spatial metric; the attempt
has not silently replaced one by the other.

For Jang-type routes, nonnegative scalar curvature after a deformation cannot
be assumed merely because a divergence integrates formally to zero. Jaracz
constructs spherical data obstructing the specific generalized-Jang/zero-
divergence system with the desired asymptotics. That result constrains that
coupling, not every Jang formulation or the Penrose inequality itself.
\eref{30006689}{2}. Boundary flux, horizon area comparison, asymptotic mass, and
rigidity must all survive the same limit. None is supplied here in general.

\Needspace{8\baselineskip}
\paragraph{Outcome.}
\textbf{Outcome: Exploratory approach with unresolved gap.}

The attempt supplies an exact mass derivative and a sharp least-enclosing-area
proof in a stated spherical untrapped subclass, plus a rigorous counterexample
to an endpoint analytic estimate used in a source discussion. Neither is a
counterexample to the Penrose inequality. A general construction of a suitable
monotone deformation or foliation, with all flux, area, and rigidity controls,
remains unproved; the requested unrestricted theorem has not been obtained.

% END RESEARCH 37

\subsection{Sources}
\begin{itemize}\raggedright
\item[\textnormal{[S1]}]\hypertarget{source:30006689:1}{} D. Xu, 'The Spacetime Penrose Inequality: Conditional Results for Stable MOTS and General Trapped Surfaces,' arXiv:2512.04137v5 (2025).
\url{https://arxiv.org/abs/2512.04137}.
{\small\textit{Catalogue formulation source.}}
\item[\textnormal{[S2]}]\hypertarget{source:30006689:2}{} The spacetime Penrose inequality under a quasi final state hypothesis.
\url{https://arxiv.org/abs/2605.18730}.
{\small\textit{Catalogue research/status reference; not a proof endorsement.}}
\item[\textnormal{[S3]}]\hypertarget{source:30006689:3}{} Proof of the Spacetime Penrose Inequality With Suboptimal Constant in the Asymptotically Flat and Asymptotically Hyperboloidal Regimes.
\url{https://arxiv.org/abs/2504.10641}.
{\small\textit{Catalogue research/status reference; not a proof endorsement.}}
% BEGIN ADDED REFERENCES 37
\item[\textnormal{[E1]}]\hypertarget{extra:30006689:1}{} Hubert L. Bray and Marcus A. Khuri,
\emph{A Jang equation approach to the Penrose inequality}, Discrete and Continuous
Dynamical Systems 27 (2010), 741--766, spherical existence theorem and proposed nonspherical coupling.
\url{https://arxiv.org/abs/0910.4785}.
\item[\textnormal{[E2]}]\hypertarget{extra:30006689:2}{} Jaroslaw S. Jaracz,
\emph{Nonexistence of Solutions to the Coupled Generalized Jang Equation/Zero Divergence System}, Classical and Quantum Gravity 40 (2023), 195013.
\url{https://arxiv.org/abs/2304.09332}.
{\small\textit{Consulted 27 September 2026; only the specified coupled system is obstructed.}}

% END ADDED REFERENCES 37
\end{itemize}

% END PRESERVED SECTION
\end{document}
